This section analyzes eight scientific works relevant to the project context. Topics include home energy management and technology adoption by users, measurement and disaggregation of electric loads, and low cost implementations with an ESP32 microcontroller.

\subsection*{A1. Revisión de aplicaciones de hogar inteligente basadas en IoT}
Alaa y colaboradores~\cite{alaa2017} realizan una revisión sistemática del panorama del hogar inteligente basado en IoT, construyendo una taxonomía de aplicaciones y señalando que la investigación se encuentra dispersa y con desafíos abiertos en interoperabilidad, costo y experiencia de usuario. El trabajo fundamenta la pertinencia de soluciones integradas y de bajo costo como la propuesta, y aporta un marco para clasificar la aplicación (monitoreo + control + gestión).

\subsection*{A2. Revisión de IoT para el hogar inteligente: retos y soluciones}
Stojkoska y Trivodaliev~\cite{stojkoska2017} revisan las tecnologías habilitantes del hogar inteligente y sus retos ---conectividad, seguridad, consumo energético de los nodos y heterogeneidad de dispositivos---, proponiendo que la arquitectura y la selección de protocolos son decisiones críticas. De este trabajo se desprende la conveniencia de usar un microcontrolador con Wi-Fi integrado y un protocolo ligero como MQTT.

\subsection*{A3. Conservación de energía mediante hogares inteligentes}
Bhati, Hansen y Chan~\cite{bhati2017} estudian el impacto de las tecnologías de hogar inteligente en la conservación de energía en hogares de Singapur, mostrando que la retroalimentación (\emph{feedback}) y la automatización contribuyen al ahorro, y que el efecto se potencia con la integración entre dispositivos, hogar y red eléctrica. Este antecedente sustenta la hipótesis de ahorro del proyecto y la importancia de visualizar la información al usuario.

\subsection*{A4. Categorías y funcionalidad de la tecnología de hogar inteligente para la gestión energética}
Ford, Pritoni, Sanguinetti y Karlin~\cite{ford2017} analizan 308 productos comerciales de gestión energética del hogar y clasifican su funcionalidad, evidenciando la brecha entre las capacidades anunciadas y las realmente útiles. El trabajo justifica priorizar un conjunto acotado de funciones bien implementadas (monitoreo, control remoto, reglas automáticas) por encima de la cantidad de prestaciones.

\subsection*{A5. Revisión sistemática de la literatura de hogar inteligente: la perspectiva del usuario}
Marikyan, Papagiannidis y Alamanos~\cite{marikyan2019} revisan la literatura desde la perspectiva del usuario, identificando factores como utilidad percibida, confianza, privacidad y facilidad de uso como determinantes de la adopción. Este antecedente orienta las decisiones de diseño de la interfaz y del asistente virtual, que deben ofrecer interacción simple y comprensible.

\subsection*{A6. Enfoques de monitoreo no intrusivo de cargas}
Zoha, Gluhak, Imran y Rajasegarar~\cite{zoha2012} presentan una revisión de las técnicas de monitoreo no intrusivo de cargas (NILM), describiendo cómo la corriente y la potencia permiten identificar el estado y la operación de los equipos. El proyecto aplica esta idea a una carga única y bien caracterizada (la terma), usando su firma electrica para inferir calentamiento y detectar anomalías.

\subsection*{A7. Gestión de energía eléctrica en el hogar inteligente}
Zipperer y colaboradores~\cite{zipperer2013} analizan las tecnologías habilitantes y el comportamiento del consumidor en la gestión energética del hogar, concluyendo que el éxito depende tanto del componente técnico como de la disposición del usuario a interactuar con la información. Refuerza la necesidad de soluciones que combinen medición, control y una interfaz amigable.

\subsection*{A8. Sistema de monitoreo de energía inteligente basado en ESP32}
El-Khozondar y colaboradores~\cite{elkhozondar2024} desarrollan un sistema de monitoreo energético de bajo costo basado en el microcontrolador ESP32, integrando sensores de tensión (ZMPT101B) y de corriente (SCT-013), y reportando los datos al usuario a través de la red, demostrando la viabilidad del ESP32 y de los transformadores de corriente para el monitoreo eléctrico doméstico. Este trabajo es el antecedente más cercano al componente de medición eléctrica del proyecto.

\subsection*{Síntesis comparativa y vacío identificado}
La Tabla~\ref{tab:arte} resume los trabajos analizados. Se observa que la literatura aborda, por un lado, la gestión energética y la adopción del usuario (A1--A5) y, por otro, la medición de cargas y las implementaciones de bajo costo (A6--A8). Sin embargo, existe un vacío en soluciones \emph{específicas para la terma doméstica} que integren simultáneamente: tres variables físicas (temperatura, corriente y flujo), un actuador de potencia seguro (SSR), comunicación inalámbrica con procesamiento previo y una interfaz que combine dashboard y asistente virtual. El presente proyecto se posiciona en ese vacío.

\begin{table*}[htbp]
    \centering
    \caption{Comparative synthesis of the state of the art.}
    \label{tab:arte}
    \footnotesize

    \begin{tabularx}{\textwidth}{@{}P{0.7cm} P{3.0cm} X P{3.0cm}@{}}
        \toprule
        \tblhead{\#} &
        \tblhead{Work} &
        \tblhead{Main contribution} &
        \tblhead{Limitation / gap} \\
        \midrule

        A1 & Alaa et al. (2017)
        & Taxonomy of IoT smart home applications
        & Not focused on a specific load \\

        A2 & Stojkoska and Trivodaliev (2017)
        & Connectivity and protocol challenges
        & No load control implementation \\

        A3 & Bhati et al. (2017)
        & Evidence of energy savings through feedback and automation
        & Dependent on household context \\

        A4 & Ford et al. (2017)
        & Functional categories of HEM (308 products)
        & No prototype development \\

        A5 & Marikyan et al. (2019)
        & User adoption factors
        & Review, not a technical system \\

        A6 & Zoha et al. (2012)
        & NILM methods for identifying loads
        & High computational complexity \\

        A7 & Zipperer et al. (2013)
        & Enabling technologies and consumer behavior
        & No low-cost prototype \\

        A8 & El-Khozondar et al. (2024)
        & ESP32-based electrical monitoring with SCT-013
        & Electrical monitoring only \\

        \bottomrule
    \end{tabularx}
\end{table*}
